\section*{\inTextCircled[10]{6}{15} Broken Span}\stepcounter{section}\phantomsection\addcontentsline{toc}{section}{Broken Span}
{\entryfont A vast section of the surrounding rock has collapsed away, splitting the tunnel across a deep subterranean chasm. The bridge that once carried the road across it has fared little better.}
\begin{DndReadAloud}
	The transport road ends abruptly at the shattered remains of a stone bridge.

	Beyond its broken edge yawns a vast chasm, its walls disappearing into darkness far below. Fractured pillars and pieces of ancient masonry cling to the rock face, while the opposite half of the bridge still reaches out from the far side.

	Between the two ends lies a wide stretch of empty air.

	Somewhere far below, loose stones strike unseen surfaces before their echoes finally fade.
\end{DndReadAloud}
\subsection*{Crossing the Chasm}\stepcounter{subsection}\phantomsection\addcontentsline{toc}{subsection}{Crossing the Chasm}
{\entryfont The gap between the surviving sections of bridge is approximately 30 feet wide. The remaining stonework is stable enough to support the expedition, though its broken edges offer little protection against a fall.

There is no single intended solution. Characters should be encouraged to use their individual abilities, equipment, spells, or the remaining bridge structure to devise a way across. Possible approaches include:
\begin{itemize}
	\item \textbf{B.A.S.A.L.T-5} Its limited flight allows it to cross the gap directly and potentially carry or assist lighter members of the expedition.
	\item \textbf{B.R.E.A.C.H-4} The Drill Arm can create secure anchor points in the surrounding stone for ropes or improvised lines.
	\item \textbf{F.I.X-2} Ropes, tools, and available bridge debris can be fashioned into a crude crossing or secured line.
	\item \textbf{L.O.R.E-3} Utility magic may allow individual characters or equipment to cross without relying on the remains of the bridge.
	\item \textbf{Silas} His light frame and agility make traversing ropes, narrow ledges, or improvised crossings considerably easier than for heavier constructs.
\end{itemize}
\noindent Other reasonable solutions should work without requiring a specific combination of characters.}